\begin{theorem}[Wild quadratic low and boundary upper-coordinate transport]
\label{mod:cases-wildquadratic-lowcoordinatetransport}
In the low and boundary ranges, transport the actual upstairs stationary
coordinate through the degree-two norm polynomial.  The construction retains
the selected low representative pair, low conductor table, and actual
upstairs witness.  It proves the correction comparison and a
\(q\)-polymorphic upper-coordinate certificate while the actual residual
coordinate is live.

At the boundary, \(\rho\) is the residue of the retained selected source
coordinate.  It is not reconstructed from the norm identity alone.  The
representatives remain at \texttt{lowCriticalFloorDepth t} and the actual
stationary depth; no stronger representative is selected.
\LeanFile{LanglandsFirstMainLemma/Cases/WildQuadratic/LowCoordinateTransport.lean}
\lean{LanglandsFirstMainLemma.wildQuadraticLowCoordinateTransport}
\leanok
\SourceText{Lemma lem:residual-parameter-transport and Propositions prop:quadratic-critical-functions and prop:quadratic-gamma-table, low and boundary upper rows.}
\uses{mod:parameters-phasereduction,mod:parameters-low,mod:cases-wildquadratic-coordinatetransport}
\FormalizationContract
Retain the actual low table, representative pair, upstairs witness, and
equalities identifying the selected correction units.  Produce the upper
coordinate and its actual extension-row phase uniformly for every refinement
\(q\).  Do not select \(q\), strengthen representative precision, or infer the
boundary source from \(N(u)=n\) alone.
\end{theorem}
